% =============================================================================
%  PROTOCOLO DE AUDITORIA MULTIPUNTO — manual de ejecucion en Windows
%  Proyecto: "Two barriers to equitable access" (UAIS)
%  Version 1.0 — 15 de agosto de 2026
% =============================================================================
\documentclass[11pt,a4paper]{article}

\usepackage[utf8]{inputenc}
\usepackage[T1]{fontenc}
\usepackage[spanish,es-noshorthands]{babel}
\usepackage{lmodern}
\usepackage{microtype}
\usepackage[a4paper,top=2.6cm,bottom=2.6cm,left=2.7cm,right=2.7cm]{geometry}
\usepackage{xcolor}
\usepackage{graphicx}
\usepackage{booktabs}
\usepackage{longtable}
\usepackage{array}
\usepackage{ragged2e}
\usepackage{enumitem}
\usepackage{titlesec}
\usepackage{fancyhdr}
\usepackage[most]{tcolorbox}
\usepackage{hyperref}

\definecolor{azul}{HTML}{1F3B73}
\definecolor{ocre}{HTML}{9E3D00}
\definecolor{grisc}{HTML}{F2F2F2}
\definecolor{grist}{HTML}{3A3A3A}

\hypersetup{
	colorlinks = true,
	linkcolor  = azul,
	urlcolor   = azul,
	pdftitle   = {Protocolo de auditoria multipunto en Windows},
	pdfauthor  = {Proyecto UAIS — privacidad y accesibilidad en sitios web universitarios}
}

\titleformat{\section}{\Large\bfseries\color{azul}}{\thesection.}{0.6em}{}
\titleformat{\subsection}{\large\bfseries\color{grist}}{\thesubsection.}{0.6em}{}
\titlespacing*{\section}{0pt}{2.2ex plus 1ex minus .2ex}{1.4ex plus .2ex}

\pagestyle{fancy}
\fancyhf{}
\fancyhead[L]{\footnotesize\color{grist}Auditoría multipunto · Protocolo de ejecución en Windows}
\fancyhead[R]{\footnotesize\color{grist}\thepage}
\renewcommand{\headrulewidth}{0.4pt}

\setlist[enumerate]{itemsep=2pt,topsep=4pt}
\setlist[itemize]{itemsep=2pt,topsep=4pt}

% --- Caja de ordenes que se teclean ----------------------------------------
\newtcblisting{orden}{
	listing only,
	colback = grisc,
	colframe = grist!35,
	boxrule = 0.4pt,
	arc = 1.5pt,
	left = 5pt, right = 5pt, top = 4pt, bottom = 4pt,
	listing options = {basicstyle=\ttfamily\small, columns=fullflexible, breaklines=true}
}

% --- Cajas de aviso ---------------------------------------------------------
\newtcolorbox{aviso}[1][]{
	colback = ocre!6, colframe = ocre, boxrule = 0.9pt, arc = 2pt,
	left = 7pt, right = 7pt, top = 5pt, bottom = 5pt,
	fonttitle = \bfseries\color{white}, coltitle = white,
	title = {ATENCIÓN #1}
}

\newtcolorbox{nota}[1][]{
	colback = azul!5, colframe = azul, boxrule = 0.9pt, arc = 2pt,
	left = 7pt, right = 7pt, top = 5pt, bottom = 5pt,
	fonttitle = \bfseries\color{white}, coltitle = white,
	title = {#1}
}

\begin{document}

\begin{titlepage}
	\centering
	\vspace*{2.5cm}
	{\Large\bfseries\color{azul} Auditoría multipunto de privacidad y accesibilidad\par}
	\vspace{0.6cm}
	{\LARGE\bfseries Protocolo de ejecución en Windows\par}
	\vspace{1.2cm}
	{\large Cómo repetir la auditoría desde Ecuador, la Unión Europea\\ y los Estados Unidos usando una VPN\par}
	\vspace{2.0cm}
	{\large Proyecto: \textit{Two barriers to equitable access}\\ Universal Access in the Information Society\par}
	\vspace{1.0cm}
	{Universidad Técnica Estatal de Quevedo\par}
	\vspace{2.5cm}
	{\normalsize Versión 1.0 — 15 de agosto de 2026\par}
	\vfill
	\begin{nota}[Para quién es este documento]
		Está escrito para que alguien sin experiencia previa en línea de órdenes pueda ejecutar la auditoría de principio a fin en un equipo con Windows y nada más. Cada orden se teclea tal cual aparece. Si algo no coincide con lo que ve en pantalla, deténgase y consulte la sección~\ref{sec:problemas} antes de continuar.
	\end{nota}
\end{titlepage}

\tableofcontents
\newpage

% =============================================================================
\section{Qué se va a medir y por qué}\label{sec:que}
% =============================================================================

La auditoría de agosto de 2026 se hizo desde un solo lugar, el Ecuador. Eso deja una pregunta sin responder: cuando una universidad europea no muestra ningún aviso de cookies a un visitante ecuatoriano y le instala rastreo sin permiso, ¿es porque trata igual a todo el mundo, o porque le muestra una cosa al visitante europeo y otra distinta al que viene de fuera?

Las dos respuestas son muy diferentes. La primera es un dato más. La segunda es un hallazgo publicable: significaría que el cumplimiento está \textbf{segmentado geográficamente}, y que la población más afectada es justamente la de aspirantes internacionales que esas universidades captan de forma activa.

Para distinguirlas hay que medir \textbf{los mismos sitios} desde \textbf{tres lugares distintos} y comparar sitio contra sitio. Eso es todo lo que hace este protocolo.

\begin{nota}[Por qué la comparación es sitio contra sitio]
	Cuando se mide el mismo sujeto en dos condiciones, los datos están \emph{pareados} (\textit{paired data}). La prueba estadística correcta entonces no es la que se usó en el artículo, sino la \textbf{prueba de McNemar}, que solo mira los casos que \emph{cambiaron} de una condición a otra. En la práctica, el resultado sale de contar cuántos sitios rastreaban visto desde Ecuador y dejaban de rastrear visto desde Europa. El programa de análisis lo calcula solo.
\end{nota}

Se harán \textbf{nueve ejecuciones} en total: tres puntos de observación (Ecuador, Unión Europea, Estados Unidos) por tres pasadas repetidas cada uno. Las tres pasadas sirven para responder a otra crítica del evaluador, la de haber medido una sola vez: con tres pasadas se puede informar cuán estable es la medida.

Cada ejecución tarda entre 30 y 50 minutos. El trabajo total son unas seis horas repartidas en tres días.

% =============================================================================
\section{Lo que hace falta antes de empezar}\label{sec:requisitos}
% =============================================================================

\begin{enumerate}
	\item Un equipo con \textbf{Windows 10 o 11}, con permisos de administrador para instalar programas.
	\item Conexión a internet estable \textbf{en el Ecuador}. Esta misma conexión será el punto de observación ecuatoriano, así que debe ser una conexión residencial o institucional normal, no una VPN.
	\item Unos 3 GB libres en disco.
	\item Una \textbf{suscripción de VPN} de pago (sección~\ref{sec:vpn}). Coste aproximado: 5 euros al mes. Basta un mes.
	\item El archivo \texttt{universidades.json} con los 126 sitios, el mismo que se usó en agosto de 2026.
	\item Los tres archivos que acompañan a este protocolo: \texttt{auditar\_\allowbreak multipunto.js}, \texttt{analizar\_\allowbreak multipunto.py} y \texttt{ejecutar.bat}.
\end{enumerate}

\begin{aviso}
	El equipo debe quedar encendido y sin usarse durante cada ejecución. Si alguien navega, ve vídeo o descarga archivos mientras corre la auditoría, los tiempos de carga cambian y algunos sitios pueden fallar. Reserve el equipo.
\end{aviso}

% =============================================================================
\section{Preparación del equipo}\label{sec:preparar}
% =============================================================================

Todo lo de esta sección se hace \textbf{una sola vez}.

\subsection{Abrir PowerShell}

Pulse la tecla \textsf{Windows}, escriba \textsf{powershell}, y en el resultado \textsf{Windows PowerShell} haga clic con el botón derecho y elija \textsf{Ejecutar como administrador}. Acepte el aviso de control de cuentas de usuario.

Verá una ventana azul o negra con un texto que termina en \texttt{>}. Ahí se teclean las órdenes. Cada orden se escribe y se pulsa \textsf{Entrar}.

\subsection{Permitir la ejecución de guiones}\label{sub:scripts}

Windows bloquea por defecto los guiones que necesita \texttt{npm}. Esto se arregla una sola vez:

\begin{orden}
Set-ExecutionPolicy -Scope CurrentUser RemoteSigned
\end{orden}

Cuando pregunte, escriba \texttt{S} y pulse \textsf{Entrar}.

\begin{aviso}[muy común]
	Si más adelante aparece el mensaje \textit{«no se puede cargar el archivo \dots\ porque la ejecución de scripts está deshabilitada en este sistema»}, es porque este paso no se hizo. Vuelva aquí.
\end{aviso}

\subsection{Instalar Node.js}

Node.js es el programa que ejecuta el auditor. Instálelo con:

\begin{orden}
winget install --id OpenJS.NodeJS.LTS -e
\end{orden}

\textbf{Cierre PowerShell y vuelva a abrirlo como administrador.} Este paso es imprescindible: hasta que no se reabre la ventana, Windows no encuentra el programa recién instalado.

Compruebe que quedó bien:

\begin{orden}
node --version
npm --version
\end{orden}

Debe responder algo como \texttt{v22.x.x} y \texttt{10.x.x}. Si dice que \texttt{node} no se reconoce, reinicie el equipo y vuelva a probar.

\begin{nota}[Si \texttt{winget} no existe]
	En equipos con Windows 10 antiguo puede no estar disponible. En ese caso descargue el instalador desde \url{https://nodejs.org} (opción \textsf{LTS}), ejecútelo y acepte todas las opciones por defecto.
\end{nota}

\subsection{Instalar Python}

Python es el que hace el análisis estadístico al final.

\begin{orden}
winget install --id Python.Python.3.12 -e
\end{orden}

Cierre y reabra PowerShell otra vez. Compruebe:

\begin{orden}
python --version
\end{orden}

Debe responder \texttt{Python 3.12.x}. Instale las dos bibliotecas que hacen falta:

\begin{orden}
python -m pip install --upgrade pip
python -m pip install scipy numpy
\end{orden}

\subsection{Crear la carpeta de trabajo}

\begin{orden}
mkdir C:\auditoria
cd C:\auditoria
\end{orden}

\begin{aviso}
	Use exactamente esa ruta. Evite carpetas dentro de OneDrive, del Escritorio o de \textsf{Documentos}: la sincronización de OneDrive puede bloquear archivos mientras el programa escribe en ellos y corromper una pasada de 40 minutos.
\end{aviso}

Copie dentro de \texttt{C:\textbackslash auditoria} estos cuatro archivos:

\begin{itemize}
	\item \texttt{auditar\_multipunto.js}
	\item \texttt{analizar\_multipunto.py}
	\item \texttt{ejecutar.bat}
	\item \texttt{universidades.json}
\end{itemize}

\subsection{Instalar el navegador de auditoría}

Estando dentro de \texttt{C:\textbackslash auditoria}:

\begin{orden}
npm init -y
npm install playwright axe-core
npx playwright install chromium
\end{orden}

La última orden descarga un navegador Chromium propio, de unos 150 MB. Es normal que tarde unos minutos.

\begin{nota}[Por qué un navegador aparte]
	El auditor no usa el Chrome que ya tiene instalado, sino una copia limpia y sin extensiones. Así la medición no queda contaminada por bloqueadores de anuncios, sesiones abiertas o cookies previas, y cualquier persona puede reproducirla exactamente.
\end{nota}

\subsection{Impedir que el equipo se suspenda}\label{sub:suspension}

Una ejecución dura 40 minutos sin que nadie toque el teclado. Si el equipo se suspende, la pasada se pierde.

\begin{orden}
powercfg /change standby-timeout-ac 0
powercfg /change hibernate-timeout-ac 0
powercfg /change monitor-timeout-ac 0
\end{orden}

Si trabaja con un portátil sin enchufar, repita cambiando \texttt{-ac} por \texttt{-dc}. Mejor aún: mantenga el equipo enchufado.

Al terminar todo el estudio puede devolver los valores originales, por ejemplo \texttt{powercfg /change standby-timeout-ac 30}.

\subsection{Prueba piloto}

Antes de gastar seis horas, compruebe que todo funciona con los cinco primeros sitios de la lista y sin VPN:

\begin{orden}
cd C:\auditoria
node auditar_multipunto.js --vantage=EC --run=0 --limit=5
\end{orden}

Debe ver, para cada sitio, una línea que empieza por \texttt{OK} y termina indicando cuántas cookies y cuántas violaciones WCAG se encontraron. Al final se generan tres archivos: \texttt{resultados\_EC\_r0.json}, \texttt{resultados\_EC\_r0.csv} y \texttt{meta\_EC\_r0.json}.

Abra \texttt{meta\_EC\_r0.json} con el Bloc de notas. Debe decir \texttt{"coincide\_con\_vantage": true} y mostrar un país \texttt{EC}. Si dice \texttt{false}, revise que no tenga una VPN encendida.

\begin{aviso}
	Cuando la prueba piloto haya salido bien, \textbf{borre} los archivos que empiezan por \texttt{resultados\_EC\_r0} y \texttt{meta\_EC\_r0}. La pasada 0 es solo un ensayo y no debe entrar en el análisis.
\end{aviso}

% =============================================================================
\section{Configuración de la VPN}\label{sec:vpn}
% =============================================================================

\subsection{Qué es y para qué sirve aquí}

Una VPN (red privada virtual) hace que todo el tráfico del equipo salga a internet por un servidor situado en otro país. Los sitios web ven la dirección de ese servidor, no la suya, y por tanto creen que usted está en Alemania o en los Estados Unidos.

Eso es exactamente lo que hace falta, porque las plataformas de consentimiento deciden si muestran o no el aviso de cookies mirando el país de la dirección IP del visitante.

\subsection{Cuál usar}

\begin{longtable}{@{}>{\RaggedRight}p{3.0cm}>{\RaggedRight}p{5.6cm}>{\RaggedRight\arraybackslash}p{5.6cm}@{}}
	\toprule
	\textbf{Opción} & \textbf{A favor} & \textbf{En contra} \\
	\midrule
	\endhead
	\textbf{La que ya tenga} \newline (McAfee, Norton, Avast\ldots) & Coste cero. Suelen permitir elegir país y son de sistema, así que sirven. & Compruebe las tres cosas del apartado siguiente. Muchas \textbf{no traen cortafuegos de emergencia}: si la VPN se cae a mitad de la pasada, el equipo sigue midiendo desde el Ecuador. \\
	\addlinespace
	\textbf{Proton VPN} \newline (plan gratuito) & Gratis y sin límite de datos. Incluye \textbf{Países Bajos}, que es Unión Europea y sirve como punto europeo, y \textbf{Estados Unidos}: exactamente los dos que hacen falta. & Más lenta. La elección de servidor es limitada. No ofrece Alemania. \\
	\addlinespace
	\textbf{Mullvad} \newline (la de mayor control) \newline unos 5 €/mes & No pide correo ni datos personales, solo un número de cuenta. Publica la lista completa de servidores. Aplicación sencilla para Windows. Se paga mes a mes y se cancela sin trámite. & Hay que pagar. \\
	\addlinespace
	\textbf{Cualquier VPN «gratis» de la tienda de aplicaciones} & --- & \textbf{No usar.} Muchas inyectan publicidad, modifican el tráfico o revenden la conexión. Contaminarían la medición. \\
	\addlinespace
	\textbf{Extensiones de VPN para Chrome} & --- & \textbf{No sirven.} Solo afectan al Chrome que usted abre, no al navegador que usa el auditor. \\
	\bottomrule
\end{longtable}

\subsection{Instalación y ajustes obligatorios}

Si ya tiene una VPN instalada, no hace falta que instale otra: sirve cualquiera que cumpla las tres condiciones del recuadro siguiente. Si no tiene ninguna, descargue la aplicación de escritorio para Windows desde el sitio oficial del proveedor que elija (\url{https://protonvpn.com/download-windows} o \url{https://mullvad.net/es/download/vpn/windows}).

\begin{nota}[Las tres condiciones que debe cumplir su VPN]
	\begin{enumerate}
		\item \textbf{Que sea de sistema, no una extensión del navegador.} Debe ser una aplicación de escritorio que desvíe todo el tráfico del equipo. Las extensiones de Chrome no afectan al navegador que usa el auditor y no sirven.
		\item \textbf{Que permita elegir país a mano.} Necesita al menos un país de la Unión Europea y los Estados Unidos. Ábrala y compruebe la lista antes de comprometerse.
		\item \textbf{Que se pueda desactivar su bloqueador de rastreadores.} Si lo trae y no se puede apagar, esa VPN no sirve para este estudio.
	\end{enumerate}
\end{nota}

Después, y antes de medir nada, revise estos cuatro ajustes. Los dos primeros son críticos.

\begin{enumerate}
	\item \textbf{Desactive el bloqueador de anuncios y rastreadores de la VPN.} En Mullvad está en \textsf{Configuración $\rightarrow$ VPN $\rightarrow$ Bloqueadores de contenido DNS}: todas las casillas deben quedar \textbf{desmarcadas}. En Proton VPN se llama \textsf{NetShield} y debe quedar \textbf{apagado}.
	\begin{quote}
		\small Si queda activado, la VPN bloquea Google Analytics y el píxel de Meta antes de que lleguen al navegador. Usted registraría cero rastreo en todas partes y concluiría que las universidades europeas son ejemplares, cuando en realidad fue su propia VPN la que las limpió. Es el error que arruinaría el estudio entero.
	\end{quote}

	\item \textbf{Active el cortafuegos de emergencia} (\textsf{kill switch} o, en Mullvad, \textsf{Modo de bloqueo permanente}). Si la VPN se cae en mitad de una pasada, el equipo se queda sin internet en lugar de seguir midiendo desde el Ecuador sin que usted se entere.

	\item \textbf{Desactive el túnel dividido} (\textsf{split tunneling}). Todo el tráfico debe pasar por la VPN.

	\item \textbf{Use el protocolo WireGuard} si la aplicación deja elegir. Es más rápido y más estable que las alternativas.
\end{enumerate}

\subsection{Si su VPN no tiene cortafuegos de emergencia}

Muchas VPN incluidas con un antivirus, entre ellas McAfee, no traen un cortafuegos de emergencia (\textit{kill switch}) real en Windows. Eso significa que, si la conexión con el servidor se cae a mitad de una pasada, el equipo vuelve a salir por su conexión ecuatoriana y el programa seguiría midiendo, pero desde el país equivocado y sin que nadie se entere.

El auditor compensa esa carencia por su cuenta. Además de comprobar el país al arrancar, \textbf{vuelve a comprobarlo cada 25 sitios y al terminar}. Si en cualquiera de esos controles detecta que la ubicación ha cambiado, imprime una alerta, la deja anotada en el archivo \texttt{meta\_} y aborta la pasada. En ese caso hay que reconectar la VPN y repetir la pasada completa, borrando antes sus tres archivos.

Todos los controles quedan registrados en \texttt{meta\_} bajo el campo \texttt{controles\_de\_ubicacion}, con el momento, el número de sitio y si fueron correctos. Es la evidencia de que la pasada entera se midió desde donde dice.

\begin{aviso}[para usuarios de McAfee]
	Compruebe además estas tres cosas en la aplicación de McAfee antes de empezar:
	\begin{enumerate}
		\item Que esté en modo de conexión \textbf{manual o siempre activa}, y no en «conectar solo en redes wifi no seguras». En una red de casa o de la universidad no se activaría sola.
		\item Que la \textbf{ubicación virtual} esté fijada al país que toca y no en «automática» o «servidor más rápido».
		\item Que el \textbf{túnel dividido} esté desactivado, si su versión lo ofrece.
	\end{enumerate}
	En la prueba piloto, compare el número de cookies que detecta con VPN y sin ella en un mismo sitio. Si con VPN sale sistemáticamente cero en todos, algo del paquete de McAfee está filtrando el tráfico y hay que buscar otra VPN.
\end{aviso}

\subsection{Qué servidores usar}

Use siempre el mismo servidor para el mismo punto de observación, en las tres pasadas. Anote el nombre exacto.

\begin{center}
	\begin{tabular}{@{}l l >{\RaggedRight\arraybackslash}p{7.4cm}@{}}
		\toprule
		\textbf{Punto} & \textbf{País} & \textbf{Qué elegir} \\
		\midrule
		\texttt{EC} & Ecuador & \textbf{VPN apagada.} Su propia conexión. \\
		\texttt{EU} & Unión Europea & Alemania si su VPN la ofrece; si no, Países Bajos, Francia o España. \\
		\texttt{US} & Estados Unidos & Cualquier ciudad; anote cuál. \\
		\bottomrule
	\end{tabular}
\end{center}

\begin{nota}[Qué país europeo elegir]
	Cualquiera de la Unión Europea vale, porque el Reglamento se aplica en todos por igual. Alemania es la primera opción si su VPN la ofrece, por la densidad de servidores y porque sus autoridades de protección de datos son de las más activas, de modo que es un lugar donde cabe esperar que las plataformas de consentimiento estén bien configuradas. Si su VPN no tiene Alemania, sirven igual Países Bajos, Francia o España. Lo único que importa es que sea siempre el mismo en las tres pasadas y que quede anotado.
\end{nota}

\subsection{Comprobar que la VPN funciona antes de cada pasada}

Con la VPN conectada, abra su navegador normal y entre en \url{https://ipinfo.io/json}. Debe mostrar el país que espera. Si muestra \texttt{EC} con la VPN encendida, la VPN no está funcionando.

No hace falta que confíe en eso: el propio auditor lo comprueba solo al arrancar y avisa si no coincide. Pero es mejor detectarlo antes de lanzar 40 minutos de medición.

\subsection{La limitación que hay que declarar en el artículo}

\begin{aviso}[honestidad metodológica]
	Las direcciones IP de una VPN pertenecen a centros de datos, no a hogares. Algunos sitios tratan a esas direcciones de forma distinta: pueden bloquearlas, mostrar un desafío antibot o servir una versión reducida de la página.

	Esto no invalida el estudio, pero \textbf{hay que declararlo} y hay que medirlo. Por eso el programa de análisis informa, en su sección 1, cuántos sitios fallaron en cada punto de observación. Si desde Ecuador cargan 126 de 126 y desde Alemania solo 108, esa diferencia debe aparecer en el artículo y debe discutirse, porque los 18 sitios que fallaron podrían no ser aleatorios.

	La alternativa más defendible ante un revisor exigente sería alquilar máquinas virtuales en la nube en regiones de la Unión Europea y de los Estados Unidos, pero eso exige cuenta con tarjeta y salirse de Windows. Con VPN el resultado es válido siempre que la limitación se declare.
\end{aviso}

% =============================================================================
\section{Protocolo de ejecución}\label{sec:ejecucion}
% =============================================================================

\subsection{Las cinco reglas de oro}

\begin{enumerate}
	\item \textbf{El mismo equipo y la misma red física} en las nueve ejecuciones. No mezcle un portátil con un sobremesa ni la red de la casa con la de la universidad.
	\item \textbf{La misma franja horaria} todos los días, por ejemplo entre las 09:00 y las 13:00. El tráfico de internet cambia a lo largo del día y eso afecta a los tiempos de carga.
	\item \textbf{Las tres ejecuciones de un mismo día, seguidas}, sin dejar horas entre medias. Así los tres puntos de observación ven los sitios en un estado casi idéntico.
	\item \textbf{Rote el orden} de los puntos cada día, según la tabla de la sección~\ref{sec:calendario}. Si midiera siempre Ecuador primero y Estados Unidos último, no podría distinguir el efecto del país del efecto de la hora.
	\item \textbf{No use el equipo} mientras corre una ejecución.
\end{enumerate}

\subsection{Calendario de las ejecuciones}\label{sec:calendario}

Lo recomendable son tres pasadas por punto de observación, nueve ejecuciones en total, repartidas en tres sesiones. En cada sesión se cambia el orden de los puntos para que ninguno coincida siempre con la misma franja horaria.

\begin{center}
	\begin{tabular}{@{}l l l l@{}}
		\toprule
		\textbf{Sesión} & \textbf{Primero} & \textbf{Segundo} & \textbf{Tercero} \\
		\midrule
		Día 1 & EC & EU & US \\
		Día 2 & EU & US & EC \\
		Día 3 & US & EC & EU \\
		\bottomrule
	\end{tabular}
\end{center}

En cada casilla se lanza \texttt{ejecutar.bat}, se elige el punto y se pide \textbf{una} pasada. El programa comprueba por sí mismo qué pasadas de ese punto ya existen y numera la nueva sin que usted tenga que llevar la cuenta, de modo que es imposible sobrescribir una medición anterior por teclear dos veces el mismo número.

También puede pedir varias pasadas seguidas del mismo punto: como todas usan la misma VPN, se encadenan solas con una pausa de sesenta segundos entre ellas y no hace falta estar delante. Es más cómodo, pero rotar entre días es metodológicamente mejor, porque separa el efecto del punto de observación del efecto de la hora.

\subsection{Qué hacer en cada ejecución, paso a paso}

\begin{enumerate}
	\item Conecte o desconecte la VPN según indique la tabla, y espere a que la aplicación diga \textsf{Conectado}.
	\item Abra PowerShell y sitúese en la carpeta:
	\begin{orden}
cd C:\auditoria
	\end{orden}
	\item Teclee la orden que corresponde a esa fila de la tabla y pulse \textsf{Entrar}.
	\item El programa comprueba primero desde qué país está midiendo e imprime el resultado. \textbf{Léalo.} Si dice \texttt{*** AVISO ***}, tiene quince segundos para pulsar \textsf{Ctrl+C} y abortar; corrija la VPN y vuelva a empezar.
	\item Deje el equipo trabajando. Verá una línea por sitio. Terminará con un resumen.
	\item Anote en la hoja de registro del anexo~\ref{sec:hoja} la fecha, la hora, el servidor de VPN y cuántos sitios fallaron.
\end{enumerate}

\begin{nota}[Atajo: el archivo \texttt{ejecutar.bat}]
	Si prefiere no teclear, haga doble clic en \texttt{ejecutar.bat}. Le preguntará el punto de observación y el número de pasada, le recordará el estado que debe tener la VPN y lanzará la orden por usted.
\end{nota}

\subsection{Si una ejecución se interrumpe}

El programa guarda después de cada sitio y sabe reanudar. Si se corta la luz o cierra la ventana por error, vuelva a teclear \textbf{exactamente la misma orden}: retomará donde quedó y saltará los sitios ya auditados.

Si lo que quiere es rehacer una pasada desde cero, borre antes los tres archivos de esa pasada, por ejemplo \texttt{resultados\_EU\_r2.json}, \texttt{resultados\_EU\_r2.csv} y \texttt{meta\_EU\_r2.json}.

\subsection{Al terminar las nueve}

En \texttt{C:\textbackslash auditoria} debe haber 27 archivos de resultados: nueve \texttt{.json}, nueve \texttt{.csv} y nueve \texttt{meta\_}. Cópielos a una carpeta de respaldo antes de tocar nada más.

% =============================================================================
\section{Análisis de los resultados}\label{sec:analisis}
% =============================================================================

\subsection{Ejecutar el análisis}

\begin{orden}
cd C:\auditoria
python analizar_multipunto.py
\end{orden}

Esto imprime el informe en pantalla y lo guarda en \texttt{analisis\_multipunto.txt}.

Para el resultado que interesa al artículo, restringido a las universidades europeas del grupo de referencia:

\begin{orden}
python analizar_multipunto.py --paises DE,FR,ES,NL,BE,SE,CH,GB --salida analisis_europa.txt
\end{orden}

Y para las ecuatorianas, que deberían comportarse igual desde los tres puntos y sirven de control:

\begin{orden}
python analizar_multipunto.py --grupo ecuador --salida analisis_ecuador.txt
\end{orden}

\subsection{Cómo leer el informe}

El informe tiene cinco secciones.

\textbf{Sección 1, control de calidad.} Una fila por ejecución. Si alguna dice \texttt{NO!} en la columna de ubicación, esa pasada se midió desde el país equivocado y hay que repetirla. Compare también la columna de fallos entre puntos: una diferencia grande es la limitación de la VPN que mencionaba la sección~\ref{sec:vpn}.

\textbf{Sección 2, estabilidad.} El porcentaje de sitios que dieron el mismo resultado en las tres pasadas. Por encima del 90\,\% la medida es sólida. Por debajo del 75\,\% hay que informarlo en el artículo como ruido de medición.

\textbf{Sección 3, proporciones.} El porcentaje de sitios que rastrean, con su intervalo de confianza, en cada punto de observación. Es la tabla descriptiva.

\textbf{Sección 4, McNemar.} Es el resultado central. Para cada par de puntos aparecen cuatro números:

\begin{itemize}
	\item \texttt{a}: sitios que rastrean vistos desde los dos puntos.
	\item \texttt{b}: rastrean desde el primero pero \emph{no} desde el segundo.
	\item \texttt{c}: no rastrean desde el primero pero \emph{sí} desde el segundo.
	\item \texttt{d}: no rastrean desde ninguno.
\end{itemize}

Los que importan son \texttt{b} y \texttt{c}, los \emph{discordantes}: los sitios que cambiaron de comportamiento según desde dónde se les mirara. Si al comparar \texttt{EC vs EU} sale un \texttt{b} grande y un \texttt{c} pequeño, eso significa que hay sitios que rastrean al visitante ecuatoriano y no al europeo. Eso, y no otra cosa, es la segmentación geográfica del cumplimiento.

\textbf{Sección 5, Q de Cochran.} Contrasta los tres puntos a la vez. Es un resumen; el detalle está en la sección 4.

\subsection{Qué esperar}

Si la hipótesis es correcta, el patrón será este: las universidades \textbf{europeas} mostrarán banner y no rastrearán vistas desde Alemania, y rastrearán sin banner vistas desde Ecuador y desde Estados Unidos. Las universidades \textbf{ecuatorianas} se comportarán igual desde los tres puntos, porque no tienen ninguna razón para geolocalizar.

Ese contraste entre los dos grupos es lo que hace convincente el resultado: el grupo ecuatoriano funciona como control negativo. Si también cambiara, habría que sospechar de la VPN antes que de las universidades.

Si, por el contrario, las europeas rastrean igual desde los tres puntos, el hallazgo es distinto pero igualmente publicable: significaría que no segmentan, que simplemente no cumplen, ni siquiera con su propio público.

% =============================================================================
\section{Qué escribir en el artículo}\label{sec:escribir}
% =============================================================================

\subsection{En Métodos}

Añada un párrafo con este contenido, adaptando los datos reales:

\begin{quote}
	\small
	\textit{The audit was repeated from three vantage points between [fechas] 2026: a residential connection in Ecuador, a WireGuard endpoint in Frankfurt, Germany (Mullvad, server \texttt{de-fra-wg-001}) and one in New York, United States (server \texttt{us-nyc-wg-201}). Each vantage point was measured three times on three separate days, with the order of vantage points rotated across days in a Latin-square design to avoid confounding vantage with time of day. The public IP address and its geolocation were recorded at the start of every run and are released with the data. DNS-level content blocking was disabled at the VPN client throughout. All other parameters of the protocol were identical to those described in Section~3.5. Per-site outcomes were consolidated across the three runs by majority rule, and paired comparisons between vantage points were tested with the exact McNemar test, with Cochran's Q for the three-way comparison.}
\end{quote}

\subsection{En Amenazas a la validez}

\begin{quote}
	\small
	\textit{The European and United States vantage points were reached through a commercial VPN, so their egress addresses belong to a data centre rather than to a residential network. Some sites treat such addresses differently, and [n] sites failed to load from [punto] against [n] from Ecuador; the affected institutions are listed in the released data. The direction of this bias is not known a priori.}
\end{quote}

\subsection{En Resultados}

Informe siempre, para cada comparación: el número de sitios pareados, los recuentos \texttt{b} y \texttt{c}, el valor $p$ exacto y la razón de momios pareada con su intervalo. Nunca solo el valor $p$.

% =============================================================================
\section{Problemas frecuentes}\label{sec:problemas}
% =============================================================================

\begin{longtable}{@{}>{\RaggedRight}p{6.0cm}>{\RaggedRight\arraybackslash}p{8.2cm}@{}}
	\toprule
	\textbf{Lo que ve} & \textbf{Qué hacer} \\
	\midrule
	\endhead
	\texttt{node no se reconoce como un comando} & Cierre y vuelva a abrir PowerShell. Si persiste, reinicie el equipo. Si sigue, reinstale Node.js. \\
	\addlinespace
	\texttt{la ejecución de scripts está deshabilitada} & Ejecute la orden del apartado~\ref{sub:scripts} y repita. \\
	\addlinespace
	\texttt{Executable doesn't exist at \dots} & Falta el navegador. Ejecute \texttt{npx playwright install chromium} dentro de \texttt{C:\textbackslash auditoria}. \\
	\addlinespace
	\texttt{Cannot find module 'playwright'} & No está en la carpeta correcta o no instaló las dependencias. Ejecute \texttt{cd C:\textbackslash auditoria} y luego \texttt{npm install playwright axe-core}. \\
	\addlinespace
	\texttt{ERROR: falta o es invalido -{}-vantage} & Escribió mal la orden. Debe ser exactamente \texttt{-{}-vantage=EU}, con dos guiones y sin espacios alrededor del igual. \\
	\addlinespace
	El programa avisa de que el país no coincide & La VPN no está conectada, o está conectada a otro país. Pulse \textsf{Ctrl+C}, corríjalo y repita. \\
	\addlinespace
	Muchos sitios dan \texttt{FALLO: \dots\ timeout} & La conexión está saturada o la VPN va lenta. Cambie a otro servidor del mismo país, anótelo, y repita la pasada completa. \\
	\addlinespace
	Windows pregunta si permite el acceso de Node a la red & Permita el acceso en redes privadas. \\
	\addlinespace
	El antivirus pone en cuarentena algo de la carpeta & Añada \texttt{C:\textbackslash auditoria} a las exclusiones del antivirus y repita la instalación. \\
	\addlinespace
	El equipo se suspendió a mitad & Aplique el apartado~\ref{sub:suspension} y relance la misma orden: reanudará donde quedó. \\
	\addlinespace
	\texttt{Falta scipy} & Ejecute \texttt{python -m pip install scipy numpy}. \\
	\addlinespace
	\texttt{No se encontro ningun archivo \dots} al analizar & Está ejecutando el análisis en otra carpeta. Haga \texttt{cd C:\textbackslash auditoria} primero. \\
	\bottomrule
\end{longtable}

% =============================================================================
\section{Anexo. Hoja de registro}\label{sec:hoja}
% =============================================================================

Imprima esta hoja y rellénela a mano después de cada ejecución. Es la bitácora del experimento y sirve como evidencia si un revisor pregunta.

\vspace{0.6cm}

\renewcommand{\arraystretch}{2.4}
\begin{center}
	\begin{tabular}{@{}p{0.5cm}p{1.3cm}p{0.9cm}p{1.6cm}p{2.8cm}p{1.3cm}p{2.4cm}@{}}
		\toprule
		\footnotesize\textbf{\#} & \footnotesize\textbf{Fecha} & \footnotesize\textbf{Pto.} & \footnotesize\textbf{Hora} & \footnotesize\textbf{Servidor VPN} & \footnotesize\textbf{Fallos} & \footnotesize\textbf{Notas} \\
		\midrule
		1 & & EC & & (sin VPN) & & \\ \hline
		2 & & EU & & & & \\ \hline
		3 & & US & & & & \\ \hline
		4 & & EU & & & & \\ \hline
		5 & & US & & & & \\ \hline
		6 & & EC & & (sin VPN) & & \\ \hline
		7 & & US & & & & \\ \hline
		8 & & EC & & (sin VPN) & & \\ \hline
		9 & & EU & & & & \\
		\bottomrule
	\end{tabular}
\end{center}
\renewcommand{\arraystretch}{1.0}

\vspace{1.0cm}

\begin{nota}[Antes de dar por cerrado el experimento]
	\begin{itemize}
		\item ¿Están los 27 archivos de resultados respaldados en dos lugares distintos?
		\item ¿Dice \texttt{true} el campo \texttt{coincide\_con\_vantage} en los nueve archivos \texttt{meta\_}?
		\item ¿Está la hoja de registro completa y escaneada?
		\item ¿Se depositaron los datos y los guiones en Zenodo u OSF, y se anotó el DOI?
	\end{itemize}
\end{nota}

\end{document}
